\documentclass[a0paper,english]{tumposter}

\usepackage{array}
\usepackage{tabularx}
\usepackage{lipsum}

\title{Very important science}
\subtitle{A promising new approach to beautiful \LaTeX\ posters}
\subject{Conference of stylish posters - Poster \#42}

\author[affil=1]{First Author}
\author[affil=2]{Second Author}
\author[affil=1]{Last Author}
\affil[mark=1]{\theSchoolName, \theUniversityName}
\affil[mark=2]{School of Applied \LaTeX, University of \TeX nology}

% You can configure your poster style globally, here in the preamble
% or you execute this command inside the document, directly before your
% poster content (if you want multiple posters in one document)
\TUMpostersetup{
  columns = 4,
  rows = 6,
  boxes = TUMBoxBlank,
  printaffil = true,
  showframe = true,
}


\begin{document}


% .:: Example poster with 4 columns
% ----------------------------------------------------------------------------

\TUMpostersetup{
  columns = 4,
  rows = 6,
  boxes = TUMBoxBlank,
  printaffil = true,
  showframe = true,
}
\maketitle
\begin{tcbposter}
  % footer with some contact information
  \posterbox[fontupper=\small, height=5cm,
    enlarge top initially by = .5\TUMposterboxsep
  ]{
    name=footer, column=1, span=4, above=bottom
  }{
    \textbf{\theUniversityName}\\
    This footer box could be used to place contact information.\\
    Setting a fixed height and insert additional spacing\\
    to separate the footer from the main content is possible.
  }

  % figure on the top right
  \posterbox{name=fig1, column=4, below=top}{
    \includegraphics[width=\columnwidth, height=30cm]{example-image-a}
    \captionof{figure}{This is a figure caption. The figure contains a
      placeholder to test the layout.}
  }

  % figure on the bottom left
  \posterbox{name=fig2, column=1, above=footer}{
    \includegraphics[width=\columnwidth, height=20cm]{example-image-b}
    \captionof{figure}{This is another figure placeholder.}
  }

  % main text box is a sequence over the rest of the page
  \posterbox{name=text, column=1,
    sequence = 1 between top and fig2 then
    2 between top and footer then
    3 between top and footer then
    4 between fig1 and footer
  }{
    \lipsum[1-15]
  }
\end{tcbposter}


% .:: Example poster with 3 columns
% ----------------------------------------------------------------------------

\clearpage
\TUMpostersetup{
  columns = 3,
  rows = 6,
  boxes = TUMBoxBlank,
  printaffil = true,
  showframe = false,
}
\maketitle
\begin{tcbposter}
  % footer with some contact information
  \posterbox[fontupper=\small]{name=footer, column=1, span=3, above=bottom}{
    \begin{tabularx}{\textwidth}{p{.3\textwidth}XXX}
      \textbf{\theUniversityName}\newline\theSchoolName
       &
      \textbf{First Author}\newline Some contact information\newline
      E-mail: \href{mailto:first.author@tum.de}{first.author@tum.de}
       &
      \textbf{Dr. Second Author}\newline Don't contact me!
       &
      \textbf{Prof. Dr. Last Author}\newline Arcisstr. 21, 80333 Munich (D)
      \newline E-mail: \href{mailto:last.author@tum.de}{last.author@tum.de}
      \newline \href{www.ei.tum.de}{www.ei.tum.de}
    \end{tabularx}
  }

  % figure on the top right
  \posterbox{name=fig1, column=3, below=top}{
    \includegraphics[width=\columnwidth, height=25cm]{example-image-a}
    \captionof{figure}{This is a figure caption. The figure contains a
      placeholder to test the layout.}
  }

  % main text box is a sequence over the rest of the page
  \posterbox{name=text,
    sequence = 1 between top and footer then
    2 between top and footer then
    3 between fig1 and footer
  }{
    \lipsum[1-16]
  }
\end{tcbposter}


% .:: Example poster with 2 columns
% ----------------------------------------------------------------------------

\clearpage
\TUMpostersetup{
  columns = 2,
  rows = 6,
  boxes = TUMBoxBlank,
  printaffil = false,
  showframe = false,
  fontsize = 28pt,
  headline = twoliner
}
\maketitle
\begin{tcbposter}
  % footer with some contact information
  \posterbox[fontupper=\small]{name=footer, column=1, span=2, above=bottom}{
    \begin{tabularx}{\textwidth}{p{.3\textwidth}XXX}
      \textbf{\theUniversityName}\newline\theSchoolName
       &
      \textbf{First Author}\newline Some contact information\newline
      E-mail: \href{mailto:first.author@tum.de}{first.author@tum.de}
       &
      \textbf{Dr. Second Author}\newline Don't contact me!
       &
      \textbf{Prof. Dr. Last Author}\newline Arcisstr. 21, 80333 Munich (D)
      \newline E-mail: \href{mailto:last.author@tum.de}{last.author@tum.de}
      \newline \href{www.ei.tum.de}{www.ei.tum.de}
    \end{tabularx}
  }

  % figure on the top right
  \posterbox{name=fig1, column=2, below=top}{
    \includegraphics[width=\columnwidth, height=25cm]{example-image-a}
    \captionof{figure}{This is a figure caption. The figure contains a
      placeholder to test the layout.}
  }

  % main text box is a sequence over the rest of the page
  \posterbox{name=text,
    sequence = 1 between top and footer then
    2 between fig1 and footer
  }{
    \lipsum[1-13]
  }
\end{tcbposter}


% .:: Example poster with single column
% ----------------------------------------------------------------------------

\clearpage
\TUMpostersetup{
  columns = 1,
  rows = 6,
  boxes = TUMBoxBlank,
  printaffil = false,
  showframe = false,
  fontsize = 28pt,
  headline = tumlogo,
}
\maketitle
\begin{tcbposter}
  % footer with some contact information
  \posterbox[fontupper=\small, height=5cm,
    enlarge top initially by = .5\TUMposterboxsep
  ]{
    name=footer, above=bottom
  }{
    \textbf{\theUniversityName}\\
    This footer box could be used to place contact information.\\
    Setting a fixed height and insert additional spacing\\
    to separate the footer from the main content is possible.
  }

  % main text box
  \posterbox[fontupper=\usekomafont{TUM.H2}, colback=TUMGrayLight]{
    name=text, between=top and footer
  }{
    \vspace{5cm}
    For posters without much text you should choose a larger font.

    \vspace{5cm}
    It is recommended to have good and self-explaining graphs and only short
    additional explanations on posters!
  }
\end{tcbposter}


% .:: Example poster with boxes (3 columns asymmetric)
% ----------------------------------------------------------------------------

\clearpage
\TUMpostersetup{
  columns = 3,
  rows = 6,
  boxes = TUMBox,
  printaffil = true,
  showframe = false,
  fontsize = 28pt,
  headline = threeliner,
}
\maketitle
\begin{tcbposter}
  % footer with some contact information
  \posterbox[TUMBox=TUMGray, fontupper=\small,
    enlarge top initially by = \TUMposterboxsep,
    equal height group=A
  ]{
    name=contact, column=1, span=1.5, above=bottom
  }{
    \textbf{\theUniversityName}\\
    This footer box could be used to place contact information.\\
    Setting a fixed height and insert additional spacing\\
    to separate the footer from the main content is possible.
  }

  % footer with acknowledgement
  \posterbox[TUMBox=TUMGray, fontupper=\small,
    enlarge top by = \TUMposterboxsep,
    equal height group=A
  ]{
    name=acknowledgement, column*=3, span=1.5, above=bottom
  }{
    \begin{minipage}[t]{.45\textwidth}
      \vspace{0pt}
      We would like to thank our sponsors\\
      Sponsor 1, Sponsor 2, and Sponsor 3, \\
      as well as the Institute of applied \LaTeX.
    \end{minipage}
    \hfill
    \begin{minipage}[t]{.54\textwidth}
      \vspace{0pt}
      \includegraphics[width=.3\textwidth, height=4cm]{example-image}
      \hfill
      \includegraphics[width=.3\textwidth, height=4cm]{example-image}
      \hfill
      \includegraphics[width=.3\textwidth, height=4cm]{example-image}
    \end{minipage}
  }

  % separator line between footer and content posterboxes
  \draw[TUMBlue, line width = \TUMposterrulewidth, line cap = round,
      dashed, dash pattern = on 0pt off 2\TUMposterrulewidth
    ] ([xshift=.5\TUMposterrulewidth]TCBPOSTER@contact.north west)
    -- +(\tcbposterwidth, 0);

  % Motivation box
  \posterbox[adjusted title=Motivation]{name=motivation,
    below=top, span=1, rowspan=3
  }{
    Here we motivate our work
  }

  % Results box
  \posterbox[adjusted title=Results]{name=results,
    column=1, between=motivation and contact
  }{
    \lipsum[5]
    \vspace{\baselineskip}

    \includegraphics[width=\textwidth, height=10cm]{example-image-b}
    \captionof{figure}{The figure caption.}
  }

  % Method and Model description
  \posterbox[adjusted title=Method and Model]{name=method,
    column=2, span=2, between=top and contact
  }{
    \lipsum[1]

    \begin{figure}
      \centering
      \includegraphics[width=.9\textwidth]{example-image-a}
      \caption{This is a figure caption.}
    \end{figure}

    \lipsum[2-3]
  }
\end{tcbposter}


% .:: Example poster with boxes (3 columns equal width)
% ----------------------------------------------------------------------------

\clearpage
\TUMpostersetup{
  columns = 3,
  rows = 12,
  boxes = TUMBox,
  printaffil = true,
  showframe = false,
  fontsize = 28pt,
}
\maketitle
\begin{tcbposter}
  % footer boxes and separator
  \posterbox[TUMBox=TUMGray, fontupper=\small,
    enlarge top initially by = \TUMposterboxsep,
  ]{
    name=contact, column=1, span=1.5, above=bottom, rowspan=1
  }{
    \textbf{\theUniversityName}\\
    This footer box could be used to place contact information.\\
    Setting a fixed height and insert additional spacing\\
    to separate the footer from the main content is possible.
  }
  \posterbox[TUMBox=TUMGray, fontupper=\small,
    enlarge top initially by = \TUMposterboxsep,
  ]{
    name=acknowledgement, column*=3, span=1.5, above=bottom, rowspan=1
  }{
    Acknowledgement
  }
  \draw[TUMBlue, line width = \TUMposterrulewidth, line cap = round,
      dashed, dash pattern = on 0pt off 2\TUMposterrulewidth
    ] ([xshift=.5\TUMposterrulewidth]TCBPOSTER@contact.north west)
    -- +(\tcbposterwidth, 0);

  % main content boxes
  \posterbox[adjusted title=Box 1]{name=box1, column=1, row=1, rowspan=6}{
    Content
  }
  \posterbox[adjusted title=Box 2]{name=box2, column=1,
    between=box1 and contact
  }{
    Content
  }
  \posterbox[adjusted title=Box 3]{name=box3, column=2,
    between=top and contact
  }{
    Content
  }
  \posterbox[adjusted title=Box 4]{name=box4, column=3, row=1, rowspan=4}{
    Content
  }
  \posterbox[adjusted title=Box 5]{name=box5, column=3, row=5, rowspan=4}{
    Content
  }
  \posterbox[TUMBox=TUMOrange, adjusted title=References
  ]{
    name=box1, column=3, between=box5 and contact
  }{
    \renewcommand{\section}[2]{}% don't print header
    \nocite{*}% insert all publications (even if they are not cited)
    \bibliographystyle{unsrt}
    \bibliography{sample}
  }
\end{tcbposter}


% .:: Example poster with boxes (blue background)
% ----------------------------------------------------------------------------

\clearpage
\TUMpostersetup{
  background = blue,
  columns = 3,
  rows = 12,
  boxes = TUMBoxInverse,
  printaffil = true,
  showframe = false,
  fontsize = 28pt,
}
\maketitle
\begin{tcbposter}
  % footer boxes and separator
  \posterbox[fontupper=\small, enlarge top initially by = \TUMposterboxsep]{
    name=contact, column=1, span=1.5, above=bottom, rowspan=1
  }{
    \textbf{\theUniversityName}\\
    This footer box could be used to place contact information.\\
    Setting a fixed height and insert additional spacing\\
    to separate the footer from the main content is possible.
  }
  \posterbox[fontupper=\small, enlarge top initially by = \TUMposterboxsep]{
    name=acknowledgement, column*=3, span=1.5, above=bottom, rowspan=1
  }{
    Acknowledgement
  }
  \draw[TUMWhite, line width = \TUMposterrulewidth, line cap = round,
      dashed, dash pattern = on 0pt off 2\TUMposterrulewidth
    ] ([xshift=.5\TUMposterrulewidth]TCBPOSTER@contact.north west)
    -- +(\tcbposterwidth, 0);

  % main content boxes
  \posterbox[adjusted title=Box 1]{name=box1, column=1, row=1, rowspan=6}{
    Content
  }
  \posterbox[adjusted title=Box 2]{name=box2, column=1,
    between=box1 and contact
  }{
    Content
  }
  \posterbox[adjusted title=Box 3]{name=box3, column=2,
    between=top and contact
  }{
    Content
  }
  \posterbox[adjusted title=Box 4]{name=box4, column=3, row=1, rowspan=4}{
    Content
  }
  \posterbox[adjusted title=Box 5]{name=box5, column=3, row=5, rowspan=4}{
    Content
  }
  \posterbox[TUMBoxInverse=TUMOrange, adjusted title=Box 6]{
    name=box1, column=3, between=box5 and contact
  }{
    Content
  }
\end{tcbposter}


% .:: Example poster with boxes (shaded background)
% ----------------------------------------------------------------------------

\clearpage
\TUMpostersetup{
  background = shaded,
  columns = 7,
  rows = 7,
  boxes = TUMBoxInverse,
  printaffil = false,
  showframe = false,
  fontsize = 28pt,
  headline = oneliner,
}
\maketitle
\begin{tcbposter}
  \posterbox[TUMBox=TUMGreen, adjusted title=Background]{name=background,
    column=1, span=7, row=1, rowspan=1
  }{
    Content
  }
  \posterbox[adjusted title=Box1]{name=box1,
    column=1, span=2, row=2, rowspan=3
  }{
    Content
  }
  \posterbox[adjusted title=Box2]{name=box2,
    column=1, span=2, row=5, rowspan=2
  }{
    Content
  }
  \posterbox[TUMBox=TUMBlueLighter, coltitle=TUMBlue, adjusted title=Outlook
  ]{
    name=outlook, column=1, span=7, row=7, rowspan=1
  }{
    Content
  }
  \posterbox[adjusted title=Box3]{name=box3,
    column=6, span=2, row=2, rowspan=3
  }{
    Content
  }
  \posterbox[adjusted title=Box4]{name=box4,
    column=6, span=2, row=5, rowspan=2
  }{
    Content
  }
  \posterbox[adjusted title=Box5]{name=box5,
    column=3, span=3, row=5, rowspan=2
  }{
    Content
  }
  \posterbox[adjusted title=Box6]{name=box6,
    column=3, span=3, row=4, rowspan=1
  }{
    Content
  }
  \posterbox[blankest]{name=fig1, column=3, span=3, row=2, rowspan=2}{
    \includegraphics[
      width=\tcbtextwidth, height=\tcbtextheight]{example-image}
  }
\end{tcbposter}

\end{document}
